
\subsubsection{3.1 Overview}

The analysis comprises three components: (1) a static analysis pipeline to extract structural errors, (2) an execution pipeline to extract behavioral errors, and (3) orthogonality measurement using Jaccard similarity.

\subsubsection{3.2 Static Analysis Pipeline}

We use pylint and mypy as static analysis tools. Pylint detects syntax errors, undefined variables, unused imports, and style violations. Mypy performs type checking.

\textbf{Configuration:}
\begin{itemize}
\item pylint: Error (E) and Warning (W) codes
\item mypy: \texttt{--ignore-missing-imports} flag, 30-second timeout
\end{itemize}

For each code sample, we extract the set of error categories detected.

\subsubsection{3.3 Execution Pipeline}

We use the EvalPlus test harness with extended test suites:
\begin{itemize}
\item HumanEval+: 164 problems, 80x more tests than original HumanEval
\item MBPP+: 399 problems, 35x more tests than original MBPP
\end{itemize}

\textbf{Configuration:} 5-second timeout per problem, 8 parallel workers.

\subsubsection{3.4 Orthogonality Measurement}

We measure overlap using Jaccard similarity:

$$J(S, E) = \frac{|S \cap E|}{|S \cup E|}$$

where S = set of problems with static errors, E = set of problems with execution errors. J = 0 indicates complete orthogonality (no overlap).

\subsubsection{3.5 Hypotheses}

We decompose the analysis into three sub-hypotheses:

\begin{itemize}
\item \textbf{H-E1 (Orthogonality):} Jaccard similarity < 0.3 indicates static and execution feedback detect largely non-overlapping error classes.
\item \textbf{H-M1 (Static Coverage):} Static analysis detects errors in > 60\% of test-failing code.
\item \textbf{H-M2 (Behavioral Detection):} Execution detects errors in > 40\% of static-clean code.
\end{itemize}

